\chapter{Petunjuk Menjalankan Sistem}
\label{chap:petunjuk}

Bab ini menyajikan panduan ringkas operasional untuk mengonfigurasi, menjalankan, dan memverifikasi seluruh komponen \textit{Disaster Data Aggregator Platform} pada lingkungan pengujian evaluator. Panduan lengkap dan komprehensif didokumentasikan pada berkas \texttt{README.md} di direktori utama repositori.

\section{Prasyarat dan Konfigurasi Lingkungan}
\label{sec:petunjuk_prasyarat}

Sebelum menjalankan sistem, pastikan perangkat kerja penguji telah memenuhi kebutuhan piranti lunak berikut:
\begin{enumerate}
    \item \textbf{Docker Engine \& Docker Compose:} Docker Engine minimal versi 24.0+ dan Docker Compose Plugin minimal versi v2.20+.
    \item \textbf{Go Toolchain (Opsional untuk kompilasi lokal):} Go v1.22+ (seluruh layanan telah memiliki kontainerisasi terisolasi sehingga instalasi Go lokal bersifat opsional).
    \item \textbf{Utilitas Terminal:} \texttt{make}, \texttt{curl}, \texttt{jq}, dan \texttt{bash} untuk mengeksekusi skrip otomatisasi.
    \item \textbf{Alat Pengujian Beban (Opsional):} \texttt{k6} (jika tidak terpasang di \textit{host}, target pengujian \texttt{make load-test} akan secara otomatis mengunduh dan menjalankan citra kontainer \texttt{grafana/k6}).
\end{enumerate}

\noindent\textbf{Langkah Konfigurasi Lingkungan (\textit{Environment Setup}):}
\begin{lstlisting}[language=bash, caption={Inisialisasi Variabel Lingkungan}]
# Salin template variabel lingkungan dari .env.example
cp .env.example .env

# (Opsional) Sesuaikan nilai rahasia atau port jika diperlukan
\end{lstlisting}

\section{Langkah Menjalankan Sistem}
\label{sec:petunjuk_menjalankan}

Seluruh orkestrasi kontainer (\textit{mock services}, basis data, broker Redpanda, layanan inti BNPB, dan konsumen hilir) dikelola melalui \texttt{Makefile} dan \texttt{docker-compose.yml}:

\begin{enumerate}
    \item \textbf{Menjalankan Seluruh Platform:}
    Eksekusi perintah berikut untuk membangun citra kontainer dan mengaktifkan 11 kontainer layanan di latar belakang:
    \begin{lstlisting}[language=bash]
make up
# Ekuivalen dengan: docker compose up -d --build
    \end{lstlisting}
    Tunggu sekitar 5--10 detik hingga kontainer \texttt{redpanda-init} selesai menginisialisasi partisi dan seluruh layanan berstatus \textit{healthy}.

    \item \textbf{Verifikasi Status Kesehatan Layanan (\textit{Health Checks}):}
    Pastikan seluruh titik akhir publik dan internal merespons dengan status \texttt{"UP"}:
    \begin{lstlisting}[language=bash]
curl -s http://localhost:8081/health | jq .   # mock-bmkg
curl -s http://localhost:8082/health | jq .   # mock-pvmbg
curl -s http://localhost:8084/health | jq .   # auth service
curl -s http://localhost:8080/health | jq .   # client-api gateway
docker compose exec client-api curl -s http://aggregator:8083/health | jq . # aggregator
    \end{lstlisting}

    \item \textbf{Melihat Log Layanan Secara Kolektif:}
    \begin{lstlisting}[language=bash]
make logs
    \end{lstlisting}

    \item \textbf{Menghentikan Sistem dan Membersihkan Lingkungan:}
    \begin{lstlisting}[language=bash]
make down
# Ekuivalen dengan: docker compose down -v
    \end{lstlisting}
\end{enumerate}

\section{Langkah Verifikasi dan Demonstrasi Solusi (Problem 1--5)}
\label{sec:petunjuk_verifikasi}

Repositori menyediakan target otomatisasi \texttt{make} terintegrasi untuk mereplikasi dan memvalidasi kriteria evaluasi setiap permasalahan:

\begin{table}[htbp]
\centering
\small
\caption{Matriks Verifikasi dan Target Eksekusi Demonstrasi Permasalahan}
\label{tab:runbook_matrix}
\begin{tabularx}{\textwidth}{@{}l l l X@{}}
\toprule
\textbf{Problem} & \textbf{Target Make} & \textbf{Skrip Pengujian} & \textbf{Kriteria Utama yang Diverifikasi} \\
\midrule
\textbf{Problem 1} & \texttt{make demo-p1} & \texttt{scripts/test\_schema\_change.sh} & Evolusi skema aditif (\texttt{confidence\_level}) diserap dinamis tanpa migrasi database atau \textit{redeploy}. \\
\addlinespace
\textbf{Problem 2} & \texttt{make demo-p2} & \texttt{scripts/test\_outage.sh} & Simulasi pemadaman PVMBG; respons terdegradasi (\texttt{degraded: true}); BMKG tetap normal. \\
& \texttt{make load-test} & \texttt{scripts/load\_test.js} & Uji beban 50 koneksi konstan selama 60 detik (k6) dengan tingkat keberhasilan 100\%. \\
\addlinespace
\textbf{Problem 3} & \texttt{make demo-p3} & \texttt{scripts/test\_cross\_credential.sh} & Isolasi \textit{trust domain} upstream (penolakan kredensial silang BMKG/PVMBG). \\
& & \texttt{scripts/test\_auth\_scopes.sh} & Penolakan akses \texttt{403} bagi Media pada \textit{raw fields}; rotasi penyegaran token JWT (TTL 60 detik). \\
\addlinespace
\textbf{Problem 4} & \texttt{make demo-p4} & \texttt{scripts/test\_storage\_restart.sh} & \textit{Independent container rebuild} dan \textit{restart}; isolasi kegagalan tanpa efek kaskade ke basis data. \\
\addlinespace
\textbf{Problem 5} & \texttt{make demo-p5} & \texttt{scripts/test\_consumer\_recovery.sh} & Distribusi \textit{fan-out}; akumulasi \textit{lag} saat \textit{consumer downtime}; pemulihan \textit{offset}; penambahan konsumen ke-3. \\
\bottomrule
\end{tabularx}
\end{table}

\noindent Setiap berkas log lengkap hasil demonstrasi pengujian otomatis ini juga diarsipkan pada direktori repositori \texttt{docs/evidence/p1/} hingga \texttt{docs/evidence/p5/}.
